% =====================================================================
%  GMN3 §13 开放问题 + 致谢 + 附录 A--E
%  原文：Lamination-Paper.tex 6290--6895 行；PDF p.107--114（印刷 p.106--113）
%  编号公式 27 条：(A.1)--(A.7)、(B.1)--(B.2)、(C.1)--(C.3)、(D.1)--(D.8)、(E.1)--(E.7)
%  本块含 \appendix（附录字母 A--E 由此产生）；附录 F、G 属 ch16
%  跨块引用一律硬编码原文编号，行尾留 % TODO，合稿时改 \eqref/\ref
% =====================================================================

\section{开放问题}
\label{sec:Open-Problems}

围绕带框 BPS 态及其保护自旋特征标，无论在其物理方面还是在数学方面，都有若干有趣的研究方向。我们在此列出其中一部分。

\begin{enumerate}

\item （motivic）Kontsevich--Soibelman 穿墙公式，预期不仅适用于 $d=4$、$\CN=2$ 场论的 BPS 简并度，也适用于 Calabi--Yau 流形上 II 型弦紧化的 BPS 简并度。遗憾的是，在超引力情形下，这个公式的清晰物理论证仍然付之阙如；尽管颇具讽刺意味的是，半本原穿墙公式最早恰恰是在超引力背景下导出的 \cite{Denef:2007vg}。因此，若能在超引力背景下找到本文所用线算子的某种类比物，那将会非常有意思。

\item 带框 PSC $\fro(L,\gamma;y)$ 与带框 BPS 指标应当有一个更为直接的数学定义。看来很清楚：仿照 \cite{Sen:1994yi,Gauntlett:1993sh,Weinberg:2006rq,Hitchin-L2} 的思路，用（奇异）单极子模空间上某些 Dirac 算子的零模——这些单极子与适当的矢量丛耦合——应当能够给出这样一个严格定义。这样的定义可以为检验我们的规范理论预言提供一个数学框架。特别地，用它来检验我们的强正性猜想，将会很有意思。\footnote{我们感谢 Edward Witten 就此所作的启发性讨论。} 也很可能还存在带框 BPS 指标及其保护自旋特征标的别的定义，它们更接近 Calabi--Yau 流形上的 Donaldson--Thomas 理论。\footnote{我们感谢 E. Diaconescu 提出这一可能性。}

\item $\CS$ 类理论是否都具有第 5.2 节% TODO 合稿改 \ref{subsec:Seeds-from-N=2}
的簇性质？能否用它来确定这类理论的 BPS 谱？特别地，能否用它解决 $r>1$ 时求 $A_r$ 理论 BPS 谱这一困难问题？

\item 形式叶层与物理的线算子之间是否存在一一对应？所有 $\CN=2$ 理论是否都是第 12 节% TODO 合稿改 \ref{sec:Tropic}
意义下的热带理论？

\item 我们相信，对于 $\sigma(\CQ)$ 被插入式 (6.4)% TODO 合稿改 \eqref{eq:Lvev-trace}
一事，或许还可以沿下面这条思路给出另一种物理解释。我们希望对取值于 $U(1)$ 的函数 $\psi_{\gamma} := e^{i\tilde\theta_{\gamma}}$ 所起的作用有一个理解，尤其想弄明白：为什么它是一个\emph{扭}同态 $T_u \to U(1)$——这里的「扭」指
$\psi_{\gamma}\psi_{\gamma'} = (-1)^{\langle \gamma, \gamma'\rangle}\psi_{\gamma+\gamma'}$。假设 $\IR^3$ 中的某一个方向被当作时间方向。那么我们就能把一个电荷为 $\gamma$ 的双荷子沿圆 $S^1$ 绝热地输运一圈。这个粒子所获得的 Aharonov--Bohm 相位——我们把它看作是对电 Wilson 线与磁 Wilson 线的测量——将是 $\psi_{\gamma}$。现在把乘积 $\psi_{\gamma}\psi_{\gamma'}$ 与 $\psi_{\gamma+\gamma'}$ 作比较，主要差别在于：由于这一对双荷子的电磁场的自旋自由度，$\gamma$ 粒子与 $\gamma'$ 粒子这一对的朴素 $(-1)^F$ 宇称，与束缚态的同名量相差一个因子 $(-1)^{\langle \gamma,\gamma'\rangle}$。把这条评注想透，会遇到若干微妙之处，我们不打算把它们一一理清。我们认为，倘若能沿上述思路给出扭同态的物理推导，那将会非常好。

\item 正如我们已经提到的，弄清楚 $\CM$ 上函数的非交换形变能否与形如式 (6.9)% TODO 合稿改 \eqref{eq:y-dep-trace}
的扭迹联系起来，将会很有意思。与此相关，也正如我们上面已经提到的，用作用在显式希尔伯特空间上的具体算子来理解诸如式 (3.38)% TODO 合稿改 \eqref{eq:form-gen}
这样的非交换生成函数，将会非常有意思。这必定牵涉到对空间 $\CM_{\CL}$ 的各种实截面的研究。

\item 正如我们已经指出的，关系式 (5.28)% TODO 合稿改 \eqref{eq:quant-skein-rel}
是量子 skein 关系，与 Chern--Simons 理论中为人熟知的那些关系密切类似。此外，非紧 Chern--Simons 的三维流形不变量涉及量子双对数 \cite{Dimofte:2009yn}。这两件事实表明，我们的结果与 Chern--Simons 理论之间可能存在联系。

\item 在 \cite{Zamolodchikov:1991et} 中，Alyosha Zamolodchikov 对某些 $Y$-系统的解陈述了一个递推猜想。有关的参考文献与背景见 \cite{keller-periodicity} 以及报告 \cite{volkov-talk}。引人注目的是，其中若干相关的方程与 Argyres--Douglas 理论的环关系相吻合。我们想提出：Zamolodchikov 的猜想可以借助 Argyres--Douglas 理论的离散对称性来理解并加以推广。这一提议也由 Sergio Cecotti 与 Cumrun Vafa 提出，并在 \cite{CNV} 中有更详细的展开。

\item 正如我们在引言中所指出的，关于本文工作与簇代数的关系，还有许多事情有待阐明。这里只举两个例子。第一，\cite{Lampe} 中的一些公式与 $SU(2)$ 规范理论中的叶层和线算子密切相关。另一个方向上，在 \cite{MR2567745} 中，Fock 与 Goncharov 陈述了一个很有意思的对偶猜想，它把两类不同的簇代数簇（cluster variety）配对起来，他们称这两类为 $\CX$ 型与 $\CA$ 型。他们的猜想似乎与「用热带标签来标记线算子 vev 的渐近行为」有关。我们希望另找机会给出 Fock--Goncharov 对偶猜想的物理解释。

\end{enumerate}

\section*{致谢}

我们感谢 Dan Freed、Emanuel Diaconescu、Alexander Goncharov、Anton Kapustin、Juan Maldacena、David Morrison、Hiraku Nakajima、Nikita Nekrasov、Tony Pantev、Nathan Seiberg、Yan Soibelman、Valerio Toledano Laredo 与 Edward Witten 的讨论。

我们要感谢 Simons 几何与物理中心（Simons Center for Geometry and Physics）主办了若干与本项目相关的精彩研讨会，并在这些研讨会期间给予我们款待。

GM 的工作由 DOE 在基金号 DE-FG02-96ER40959 下资助。AN 的工作部分由 NSF 在基金号 PHY-0503584 与 PHY-0804450 下资助。DG 部分由 NSF 基金 PHY-0503584 资助。DG 部分由高等研究院（Institute for Advanced Study）的 Roger Dashen 会员资格资助。

\appendix

\section{关于 $d=4$、$\CN=2$ 超对称的一些技术细节}
\label{app:Susy-Conventions}

我们采用 Bagger 与 Wess 关于 $d=4$、$\CN=1$ 超对称的约定 \cite{Wess:1992cp}。具体而言，$SU(2)$ 指标用 $\epsilon^{12}= \epsilon_{21}=1$ 升降。$so(1,3)$ 的不可约自旋表示中张量的分量记作 $\alpha,\dot\alpha$，取值 $1,2$。取复共轭的规则是 $(\CO_1 \CO_2)^\dagger = \CO_2^\dagger \CO_1^\dagger$ 以及 $(\psi_\alpha)^\dagger = \bar \psi_{\dot \alpha}$。

$\CN=2$ 超对称算子是 $(Q_\alpha^{~A}, \bar Q_{\dot \alpha B})$，其中 $A,B$ 是 $SU(2)_R$ 指标，取值 $1$ 到 $2$。它们满足厄米性条件
%
\begin{equation}
(Q_\alpha^{~A})^\dagger = \bar Q_{\dot \alpha A}
\end{equation}
%
以及 $\CN=2$ 代数
%
\begin{equation}
\begin{split}
\{ Q_{\alpha}^{~A}, \bar Q_{\dot \beta B} \} & = 2
\sigma^m_{\alpha\dot\beta}P_m \delta^{A}_{~B} \\
\{ Q_{\alpha}^{~A}, Q_{  \beta}^{~ B} \} & = 2 \epsilon_{\alpha\beta}\epsilon^{AB} \bar Z \\
\{ \bar Q_{\dot\alpha A}, \bar Q_{  \dot \beta  B} \} & = -2
\epsilon_{\dot\alpha\dot\beta}\epsilon_{AB}   Z\\
\end{split}
\end{equation}
%
其中 $Z$ 是中心荷，$P_m$ 是厄米的能动量矢量，$P^0 \geq 0$。

把线算子插入空间坐标的原点 $x^i =0$ 处，它保持一个 $so(3) \oplus su(2)_R$ 对称性，以及下列超对称：
%
\begin{equation}
\CR_{\alpha}^{~A} = \xi^{-1} Q_{\alpha}^{~A} + \xi
\sigma^0_{\alpha\dot\beta} \bar Q^{\dot \beta A}
\end{equation}
%
这里 $\xi$ 是一个相位：$\vert \xi \vert =1$。这些算子满足厄米性条件
%
\begin{equation}
\begin{split}
(\CR_{1}^{~1})^\dagger & = - \CR_{2}^{~2} \\
(\CR_{1}^{~2})^\dagger & =   \CR_{2}^{~1} \\
\end{split}
\end{equation}
%
以及代数
%
\begin{equation}
\{ \CR_{\alpha}^{~A} , \CR_{\beta}^{~B} \} = 4 \left( E + \re(Z/\zeta) \right) \epsilon_{\alpha\beta} \epsilon^{AB}
\end{equation}
%
其中 $E=P^0$ 是能量算子，而 $\zeta = \xi^{-2}$。这个代数蕴涵
%
\begin{equation}
\left( \CR_{1}^{~1} + (\CR_{1}^{~1})^\dagger \right)^2 = \left(
\CR_{1}^{~2} + (\CR_{1}^{~2})^\dagger \right)^2 =  4 (E + \re(Z/\zeta) )
\end{equation}
%
由此即得 BPS 界
%
\begin{equation}
E + \re(Z/\zeta) \geq 0.
\end{equation}

\section{线算子 vev 的全纯性}
\label{app:Holomorphy}

本附录要详细解释：被 $\CR_{\alpha}^{~A}$ 湮灭的算子，其关联函数 $\inprod{L_\zeta }$ 如何可以被视为复结构 $\zeta$ 下 $\CM$ 上的全纯函数。我们还会就它与三维理论的 Rozansky--Witten 扭转之间的关系略作评论。

做这一辨认的关键一步，是把四维超荷与低能三维 sigma 模型的超荷对上，后者的靶空间是超凯勒的 $\CM$。在圆紧化破坏四维洛伦兹不变性之后，四维超荷可以收拢成一个二重态 $\left(Q_{\alpha}^{~A}, \sigma^0_{\alpha\dot\beta} \bar Q^{\dot \beta A}\right)$，我们可以把它记作 $Q^{~A}_{a~\alpha}$，其中新指标 $a$ 取值 $1,2$。四维旋量指标 $\alpha$ 可以等同于一个三维旋量指标，而 $Q^{~A}_{a~\alpha}$ 可以等同于三维 ${\cal N}=4$ 理论的超荷。注意：$A$ 指标受四维 $SU(2)_R$ R 对称性的作用，而后者仍是三维理论的对称性；但一般而言，四维理论中并没有任何对称性把 $a=1,2$ 的超荷相互旋转。$a$ 指标\emph{不是}任何 $SU(2)$ 对称性的指标。四维 $U(1)_R$ 对称性——如果未被质量或规范耦合标度破坏的话——当然会把 $a=1,2$ 两个分量朝相反方向旋转。

这些事实与靶空间为一般超凯勒流形的三维 sigma 模型的一般性质恰好吻合。我们稍后会详细说明：这类模型总有一个旋转费米子场的 $SU(2)_R$ R 对称性，但第二个 R 对称群必须取玻色靶空间的一个等距的形式，因为它要把三个超凯勒形式相互旋转。因此，只有当超凯勒靶流形具有特殊性质时，才可能有第二个 R 对称群。它可以是一个 $U(1)_R$ 等距，把三个超凯勒形式中的两个相互旋转；也可以是一个 $SU(2)_{R}'$ 等距群，把三个超凯勒形式全都旋转。

在 sigma 模型 $\CM$ 的靶空间中，黎曼结构群 $SO(4n)$ 约化为 $(SU(2)\times USp(2n))/\IZ_2 $，而切空间 $TX\otimes\IC \cong S \otimes V$，其中 $S$ 是平凡的二秩复矢量丛，$V$ 是秩 $2n$、结构群为 $USp(2n)$ 的复矢量丛。在 $S$ 中选取一条直线，就确定了 $\CM$ 上的一个复结构。这些复结构由 $\IC P^1$ 参数化，我们让 $\zeta $ 是这个扭量球面上的非齐次坐标。局部地，我们可以引入矢量场 $W_{a i}$，其中 $i=1,\dots, 2n$ 张成 $V$，$a=1,2$ 张成 $S$。费米子场 $\psi_{\alpha}^{A i}$ 是一个 $SU(2)_R$ 二重态的时空旋量，它们是 $V$ 的截面。作用在 $\CM$ 上函数 $F$ 的超对称变换的一般形式是
%
\begin{equation}
[Q^{~A}_{a~\alpha} , F] = \psi_{\alpha}^{A i} W_{a i} F
\end{equation}
%
特别地，超荷也是平凡丛 $S$ 的截面。

把四维矢量多重态约化到三维时，我们很容易还原出这一结构。四维费米子确实按 $SU(2)_R$ 的二重态变换。$n$ 个四维费米子 $\psi_\alpha^{A}$ 与 $n$ 个复共轭的 $\sigma^0_{\alpha\dot\beta} \bar \psi^{\dot \beta A}$ 合并起来，构成辛丛 $V$ 的 $2n$ 个分量。可以验证：$Q_\alpha^A$ 在复结构 $\zeta=0$ 下作用为 $\CM$ 上函数的全纯微分算子——它显然湮灭反全纯的四维标量，而且不那么显然地，它也湮灭电 Wilson 线与磁 Wilson 线的反全纯组合。类似地，$\bar Q^{\dot \beta A}$ 在复结构 $\zeta=0$ 下作用为反全纯微分算子。到此就完成了四维超荷的指标 $a$ 与三维 sigma 模型变换律中的指标 $a$ 的辨认。这一辨认不受量子修正的影响：量子修正使靶空间保持超凯勒，并且不改变 $\zeta=0$ 处的复结构。

特别地，选取一组四维超荷 $\CR_{\alpha}^{~A}$，等同于在 $S$ 中选取一条直线，也就等同于在 $\CM$ 中选取一个复结构。式
%
\begin{equation}
[\CR_{\alpha}^{~A} , F] = \psi_{\alpha}^{A i} \left( \xi^{-1} W_{1 i}+ \xi W_{2 i} \right) F
\end{equation}
%
中出现的那些矢量场，在复结构 $\zeta = \xi^2$ 下是反全纯的。因此 $\langle L(\zeta,\dots ) \rangle$ 是复结构 $\zeta$ 下 $\CM$ 上的全纯函数。

与三维 sigma 模型沿 Rozansky--Witten 路线 \cite{Rozansky:1996bq} 所作的拓扑扭转版本作一比较，也许会有意思。拓扑扭转把三维洛伦兹群换成旧的三维洛伦兹群与 $SU(2)_R$ 的对角组合。此时有 $\IC P^1$ 那样多的拓扑荷，它们与 $\CR_{\alpha}^{~A} \delta^{\alpha}_{~A}$ 重合。显然，我们这里考虑的四维线算子会在三维理论中给出点状的拓扑可观测量。这也许会被证明是一个有用的观点。

\section{不动点方程}
\label{app:Fixed-Point-Equations}

一个 $\CN=2$ 矢量多重态含有一个标量 $\varphi$、费米子 $\psi_{\alpha A}$（处于 $so(1,3)\oplus su(2)_R$ 的 $(\mathbf{2};\mathbf{1})\otimes 2$ 之中，以及它们的复共轭 $\bar \psi_{\dot \alpha A} := (\psi_{\alpha}^{~A})^\dagger$）、一个厄米规范场 $A_m$，以及一个满足实性条件 $(D_{AB})^* = - D^{AB}$ 的辅助场 $D_{AB} = D_{BA}$。乘以 $i$ 之后，所有这些场都取值于伴随表示。\footnote{我们的约定是：$u(N)$ 的生成元是 $N\times N$ 反厄米矩阵。在协变导数为 $d +A$ 的「几何」约定中，本文的规范场与它由 $A^{\rm geometric} = i g A^{\rm hermitian}$ 相联系。} 超对称变换律为\footnote{这里我们与 Bagger--Wess 约定略有偏离。我们的标量场 $\varphi=\sqrt{2}A^{\rm BW}$，其中 $A^{\rm BW}$ 是 BW 手征多重态的标量分量。}
%
\begin{equation}
\begin{split}
[Q_{\alpha A}, \varphi] & = - 2 \psi_{\alpha A} \\
[\bar Q_{\dot\alpha A}, \varphi] & = 0 \\
[Q_{\alpha A}, A_m] & = i \bar \psi_{\dot \beta A} (\bar
\sigma_m)^{\dot \beta}_{~~\alpha} \\
[\bar Q_{\dot \alpha A}, A_m] & = -i (\bar
\sigma_m)_{\dot \alpha}^{~~\beta}  \psi_{ \beta A}  \\
 [Q_{\alpha A}, \psi_{\beta B} ] & =
\sigma^{mn}_{\beta\alpha} F_{mn} \epsilon_{AB} + i D_{AB}
\epsilon_{\beta\alpha} + \frac{i}{2} g
\epsilon_{\beta\alpha} \epsilon_{AB} [\varphi^\dagger, \varphi] \\
[\bar Q_{\dot \alpha A}, \psi_{\beta B} ] & = - i   \epsilon_{AB} \sigma^m_{\beta\dot\alpha} D_m\varphi \\
[Q_{\alpha A}, D_{BC}  ] & = \left( \epsilon_{AB}
\sigma_{\alpha}^{m\dot \beta} D_m \bar\psi_{\dot \beta C} +
B\leftrightarrow C \right) \\
& +   g \left(\epsilon_{AB} [ \varphi^\dagger, \psi_{\alpha C}] +
B\leftrightarrow C   \right)\\
\end{split}
\end{equation}

超对称 $\CR_{\alpha}^{~A}$ 的不动点方程是
%
\begin{equation}\label{eq:Q-zeta-fix}
\begin{split}
F_{0\ell} - \frac{i}{2} \epsilon_{jk\ell} F_{jk} - i   D_\ell
( \varphi/\zeta) & = 0\\
D_0(  \varphi/\zeta) - \frac{g}{2} [  \varphi^\dagger,
  \varphi ] & = 0 \\
\end{split}
\end{equation}

若结构群 $\lieg$ 约化到嘉当子代数，我们就能写出一个静态解
%
\begin{equation}\label{eq:Embed-fix-sol}
\begin{split}
F & = \frac{1}{2} \omega_S \otimes \rho^M \\
\varphi & = \zeta \frac{\rho^M}{r} + \varphi_\infty \\
\end{split}
\end{equation}
%
其中 $\rho^M\in \liet$，而 $\omega_S = \sin\theta\, d\theta\, d\phi$，这里用的是围绕原点的标准角坐标。若 $\rho^M \in \Lambda_G$，则 $A$ 会被适当地量子化。这个解定义了紫外边界条件，从而定义了一个由 $\Gamma_G$ 中某个元素标记的 't Hooft 算子。

（这些方程里会出现一个有意思的微妙之处。若取 $\rho^M, \rho^E\in \liet$ 为常数，并令 $\varphi = \zeta(\rho^M-i \rho^E)/\sqrt{2} + \varphi_\infty$ 和 $F = \half \omega_S \otimes \rho^M + \half * \omega_S \otimes \rho^E$，那么不动点方程与 Bianchi 恒等式都得到满足，但运动方程\emph{并不}满足！）

\section{低能有效理论中的不动点方程}
\label{app:Low-Energy-FxdPoint}

不动点方程 \eqref{eq:Q-zeta-fix} 在红外理论中同样成立，但其解释略有不同。

让我们回忆：红外理论中的电荷格 $\hat \Gamma$ 带有一个反对称形式 $\langle \cdot, \cdot \rangle $，并构成一个正合列
%
\begin{equation}
0 \rightarrow \Gamma_f \rightarrow \Gamma \rightarrow
\Gamma_{\rm gauge} \rightarrow 0
\end{equation}
%
其中味荷格位于 $\langle \cdot, \cdot \rangle $ 的零化子之中，而商 $\Gamma_{\rm gauge} $ 是辛的。我们把投影映射记作 $\gamma \to \bar \gamma$。

低能的 Seiberg--Witten 有效红外理论是一个自对偶的阿贝尔规范理论，其不变场强 $\CF \in \Omega^2(\IR^{1,3}) \otimes V$，其中 $V= \Gamma_{\rm gauge} \otimes \IR$ 是一个辛矢量空间。模的取值确定了 $V$ 上的一个相容复结构 $\CI$。场强满足 $d \CF =0$ 以及反自对偶约束：
%
\begin{equation}
(*\otimes \CI)\CF = - \CF
\end{equation}
%
此外，在由 $\gamma_c$ 标记的希尔伯特空间扇区中，有一个量子化条件：
%
\begin{equation}
\int_{S^2_{\infty}} \frac{\CF}{2\pi}  =  \overline{\gamma_c}.
\end{equation}

在低能有效理论中，不动点方程 \eqref{eq:Q-zeta-fix} 的一个解可以通过取
%
\begin{equation}
\CF =\frac{1}{2} \left( \omega_S \otimes \gamma_c - \omega_H \otimes
\CI(\gamma_c) \right)
\end{equation}
%
得到，其中 $\omega_H = *_4 \omega_S = {dr \wedge \frac{dt}{r^2}}$。只要矢量多重态的模按方程
%
\begin{equation}\label{eq:Mod-ATT}
2 {\rm Im} \left[ \zeta^{-1} Z(\gamma; u(r))\right] = -
\frac{\langle \gamma, \gamma_c\rangle }{r} +2 {\rm Im} \left[
\zeta^{-1} Z(\gamma; u)\right] \qquad\qquad \forall \gamma \in \hat
\Gamma
\end{equation}
%
随 $r$ 变化，不动点方程就得到满足。此式是把方程 \eqref{eq:Q-zeta-fix} 取虚部得到的。那个方程是在一个固定的对偶标架中写出的。把方程写成对偶不变的形式，就得到 \eqref{eq:Mod-ATT}。

方程 \eqref{eq:Mod-ATT} 是标准吸引子方程的一个变形。通常的吸引子方程是在 $\CN=2$ 超引力中写出的。那个方程的场论极限给出形如 \eqref{eq:Mod-ATT} 的方程，但有一个重要的例外：$\zeta = e^{i \alpha}$ 是中心荷 $Z(\gamma_c;u)$ 的（常数）相位，并且 $\Gamma_f=0$。吸引子方程的这一推广并不完全是新的。「orientihole」的方程同样把吸引子方程中的相位 $e^{i \alpha}$ 换成 $-e^{i \alpha}$ \cite{Denef:2009ja}。

一个电荷为 $\gamma_h\in \hat \Gamma $ 的探针 BPS 粒子，在上述某个场位形中运动时，其动力学由作用量
%
\begin{equation}
\int \vert Z(\gamma_h;u(r))\vert ds + \int \langle
\overline{\gamma_h}, \CA\rangle
\end{equation}
%
支配，其中积分沿探针粒子的世界线进行。因此，这样一个粒子静止时的能量是
%
\begin{equation}
\begin{split}
E  & =  \vert Z(\gamma_h;u(r))\vert - \langle \overline{\gamma_h},
\CA_0 \rangle\\
& =\vert Z(\gamma_h;u(r))\vert  + \frac{ (
\overline{\gamma_h},\overline{\gamma_c} )}{r}
 \\
\end{split}
\end{equation}
%
第二项 $\frac{ ( \overline{\gamma_h},\overline{\gamma_c}
)}{r}$ 是库仑能，它用 $V$ 上由复结构构成的正定对称度量来表达：$(v_1, v_2) := \langle v_1, \CI v_2 \rangle$。注意这个表达式对电荷是对称的、且正定，这在物理上是合理的。

现在，把不动点方程 \eqref{eq:Q-zeta-fix} 取实部，并写出其对偶不变的延拓，我们得到
%
\begin{equation}
\re\left[ \zeta^{-1} Z(\gamma_h;u(r))\right] =\frac{ (
\overline{\gamma_h},\overline{\gamma_c} )}{r} +\re \left[
\zeta^{-1} Z(\gamma_h;u)\right]
\end{equation}
%
这样我们就导出了 halo 粒子在 $\zeta$-吸引子背景中能量的公式 (3.22)% TODO 合稿改 \eqref{eq:Halo-Particle-Energy}
。halo 半径 (3.23)% TODO 合稿改 \eqref{eq:halo-rad}
则由把 \eqref{eq:Mod-ATT} 的右端对 $\gamma=\gamma_h$ 置零而得到。

\section{六维超对称}
\label{app:Six-Super}

$\IR^{1,5}$ 中的 $(2,0)$ 超共形代数（此处记作 $osp(8^*\vert 4)$）有偶子代数 $osp(8^*\vert 4)^0 = so(6,2) \oplus so(5)$。为构成奇子空间，我们选取 $so(6,2)$ 的两个手征旋量表示 $\Delta_\pm$ 之一。$\Delta_\pm$ 是四元数空间，各自同构于 $\mathbb{H}^4$。为确定起见，取 $\Delta_+$。再令 $\Delta'\cong \mathbb{H}^2$ 为 $so(5)$ 的不可约旋量表示。奇子空间是 $osp(8^*\vert 4)^1 = \Delta_+ \otimes_{\IR} \Delta'$。它是 $32$ 实维的。

庞加莱超对称的子超代数有偶子代数 $so(5,1) \oplus so(5)$；此时出现 $so(5,1)$ 的一个手征旋量 $\Delta_+'\cong \mathbb{H}^2$，用它来定义 $16$ 个庞加莱超荷构成的奇子空间 $\Delta_+'\otimes \Delta'$。

引入 Clifford 代数 $Cl(1_+,5_-)$ 的一个表示会很有用，它由形如
%
\begin{equation}
\Gamma^M = \begin{pmatrix} 0 & \bar\gamma^{M~\dot s}_{r}   \\    \gamma^{M~ s}_{\dot r} & 0 \\ \end{pmatrix}
\end{equation}
%
的 $8\times 8$ 复值矩阵给出，其中 $r=1, \dots,4 $ 标记手征自旋表示，而 $\dot r = 1, \dots, 4$ 标记反手征表示。当然 $M=0,1,\dots 5$ 跑遍时空维度。于是 $(\Gamma^M)^*  = B \Gamma^M B^{-1} $，其中
%
\begin{equation}
B = \begin{pmatrix} B_{r}^{~~s} & 0 \\ 0 & B_{\dot r}^{~~ \dot s} \\ \end{pmatrix}
\end{equation}
%
手征表示 $\Delta_\pm'$ 是伪实的。此外还有
%
\begin{equation}
C = \begin{pmatrix} 0 & c^{r \dot s} \\ \bar c^{\dot r s} & 0 \\ \end{pmatrix}
\end{equation}
%
使得 $(C \Gamma^M)^{\rm tr} = - (C \Gamma^M)$ 且 $(\Gamma^M C^{-1})^{\rm tr} = - (\Gamma^M C^{-1})$。张量 $c^{r \dot  s}$ 给出 $\Delta_+'$ 与 $\Delta_-'$ 缩并为单态的不变缩并。

对于 R 对称性 $so(5)\cong usp(4)$，把旋量 $\Delta'$ 与 $USp(4)$ 的 $4$ 等同起来。若要引入指标，我们令 $i,j =1,\dots, 4$。Gamma 矩阵的指标结构是 $\Gamma^{Ii}_{~~j}$，其中 $I=1, \dots 5$。指标用辛矩阵 $J^{ij}$ 及其逆矩阵升降。于是 $(\Gamma^I)_{i}^{~j} = J_{ii'} J^{jj'} \Gamma^{Ii'}_{~~j'}$。我们有 $(\Gamma^{Ii}_{~~j})^* = - \Gamma^{I~j}_{i}$，并且 $\Gamma^I_{ij} := J_{ii'} \Gamma^{Ii'}_{~~j}$ 是反对称的。

若把庞加莱超对称记作 $Q_{r}^i$，则我们要求
%
\begin{equation}
(Q_{r i})^\dagger = B_r^{~~s} J^{ij} Q_{s j}
\end{equation}
%
于是有 $16$ 个厄米超荷。基本的超对称代数是
%
\begin{equation}
[Q_{r}^i, Q_{s}^j ] =  2i  (\bar\gamma^M\bar c)_{rs} J^{ij} P_M
\end{equation}
%
其中 $P_M$ 是厄米的平移算子。系数 $2i$ 由约化到上面描述的 $d=4$、$\CN=2$ 代数来确定。作用在场多重态上时，右端可能因规范变换而出现额外项。作用在孤子扇区的态上时，右端可能因「中心扩张」而出现额外项。事实上，在存在弦激发时，庞加莱超代数被扩张为
%
\begin{equation}
[Q_{ri}, Q_{sj} ] =  2i  (\bar\gamma^M\bar c)_{rs} J^{ij} P_M +2  (\bar\gamma^M\bar c)_{rs} \Gamma^{I ij} Z_M^I
\end{equation}
%
其中 $Z_M^I$ 对应于 BPS 弦的「中心荷」。

超共形代数有一个场表示，称为「张量多重态」，通常记作 $(B, \psi, Y)$，其中 $B$ 是 $\IR^{1,5}$ 上一个局部定义的实二形式势，其场强是反自对偶的（定向取 $\epsilon^{012345}=+1$）；$\psi \in (\Delta_+'\otimes \Delta')$ 满足一个实性约束；$Y$ 是一个实标量场，属于 $so(5)$ 的 $5$ \cite{Howe:1983fr}。若把指标写出来，则张量多重态中的场记作 $(B_{MN}, \psi_{r}^i, Y^I)$。庞加莱超对称下的变换律为：
%
\begin{equation}
\begin{split}
[Q_{r}^i, Y^I] &=   i \Gamma^{Ii}_{~~ j} \psi_{r}^j \\
[Q_{r}^i, \psi_{sj} ]  & =   i (\bar \gamma^M\bar c)_{rs} \Gamma^{Iij} \partial_M Y^I  - \frac{i}{12} J^{ij} (\bar\gamma\bar c)^{MNP}_{rs} H_{MNP} \\
[Q_{r}^i, B_{MN} ] & =  i (\gamma_{MN})_r^{~~s} \psi_{s}^{i} \\
\end{split}
\end{equation}
%
$Y^I$ 的标度维数是 $2$，$\psi$ 的标度维数是 $5/2$，而 $B_{MN}$ 的标度维数是 $2$，于是 $B = \half B_{MN} dx^M \wedge dx^N$ 与 $H:=dB$ 都是无量纲的。场强 $H$ 必须是反自对偶的；在我们的归一化下，对于 $U(1)$ 张量多重态，$H$ 在具有非平凡拓扑的流形上有整周期。

这些超对称变换只在壳上闭合；$[Q_{ri}, Q_{sj} ]$ 作用在 $B_{MN}$ 上的闭合性只到差一个规范变换的程度，该变换由 $\Lambda_N = -2 (\bar \gamma_N\bar c)_{sr} \Gamma^{Iij} Y^I + 2 J^{ij} (\bar \gamma_P\bar c)_{sr}B_{NP}$ 给出。作用在弦态上时，所诱导的规范变换通过测量弦的电荷而起作用。这就导致中心荷 $Z_{M}^I= \hat t_M Y^I$，由此可以算出弦的张力 $\sqrt{Y^I Y^I}$。

五维的超庞加莱代数有玻色子代数 $so(1,4) \oplus so(5)$。$so(1,4)$ 的自旋表示 $\Delta''$ 同样同构于 $\mathbb{H}^2$，而奇子空间是 $\Delta''\otimes \Delta'$，对应 $16$ 个实超对称。超 Yang--Mills 多重态是 $(A,\chi,Y)$，其中 $A$ 是一个联络，$\chi \in \Delta''\otimes \Delta'$，$Y$ 是一个标量，它们都取值于一个李代数 $\lieg$。我们刻意对标量沿用了同一个记号 $Y$，因为张量多重态作维数约化会给出一个 $u(1)$ 矢量多重态。不过，在后一种情形里，更好的做法是把标量归一化为 $\Phi^I = R Y^I$，其标度维数为 $1$。
